\def\stat{mizin}





\def\tit{О ВКЛАДЕ АКАДЕМИКА И.\,А.~МИЗИНА В~ТЕОРИЮ И~ПРАКТИКУ СОЗДАНИЯ


ОТЕЧЕСТВЕННЫХ ИНФОРМАЦИОННО-ТЕЛЕКОММУНИКАЦИОННЫХ СИСТЕМ:\\


К~80-ЛЕТИЮ СО ДНЯ РОЖДЕНИЯ}





\def\titkol{О вкладе академика И.\,А.~Мизина в~теорию и~практику создания


отечественных ИТС} %: К~80-летию со дня рождения}





\def\autkol{И.\,А.~Соколов, А.\,А.~Зацаринный, В.\,Н.~Захаров}





\def\aut{И.\,А.~Соколов$^1$, А.\,А.~Зацаринный$^2$, В.\,Н.~Захаров$^3$}





\titel{\tit}{\aut}{\autkol}{\titkol}





%{\renewcommand{\thefootnote}{\fnsymbol{footnote}}


%\footnotetext[1]


%{Статья рекомендована к~публикации в~журнале Программным комитетом конференции <<Tools \& Methods


%of Program Analysis>> (<<Инструменты и~методы анализа программ>>, TMPA-2014), 14--15~ноября


%2014~года, г.~Кострома, РФ.}}





\renewcommand{\thefootnote}{\arabic{footnote}}


\footnotetext[1]{Федеральныйо исследовательский центр <<Информатика


и~управление>> Российской академии наук, isokolov@ipiran.ru}


\footnotetext[2]{Федеральный исследовательский


центр <<Информатика и~управление>> Российской академии наук,


azatsarinny@ipiran.ru}


\footnotetext[3]{Федеральный исследовательский


центр <<Информатика и~управление>> Российской академии наук,


vzakharov@ipiran.ru}





\vspace*{-6pt}





\label{st\stat}





                  \thispagestyle{headings}





\Abst{Статья посвящена 80-летнему юбилею академика И.\,А.~Мизина, директора


ИПИ РАН в~период 1989--1999~гг., выдающегося ученого, конструктора и~инженера. Дана краткая биографическая справка, характеризующая


на\-уч\-но-прак\-ти\-че\-скую деятельность И.\,А.~Мизина. Вклад академика


И.\,А.~Мизина в~теорию и~практику создания отечественных


ин\-фор\-ма\-ци\-он\-но-те\-ле\-ком\-му\-ни\-ка\-ци\-он\-ных


сетей рассматривается применительно к трем


направлениям его деятельности. Первое~--- разработка и~внедрение системы


обмена данными (СОД) в~интересах автоматизированных систем управления (АСУ)


Вооруженных Сил, первой отечественной сети


с~пакетной коммутацией. Второе направление~--- обоснование информационных


технологий для создания крупных территориальных систем передачи данных


в~регионах России. Третье направление~--- создание и~развитие информационных


сетей в~интересах органов государственной власти с~учетом требований по защите


информации.}





\KW{главный конструктор; академик;


информационно-те\-ле\-ком\-му\-ни\-ка\-ци\-он\-ные сети; информационные


технологии; система обмена данными; пакетная коммутация; методы передачи


информации; защита информации; каналы связи; испытания; органы


государственной власти}





\DOI{10.14357\!/\!08696527150114}





\vspace*{-9pt}





\vskip 10pt plus 9pt minus 6pt





      \thispagestyle{headings}








\section{Введение}





\vspace*{-3pt}








   12 апреля 2015~г.\ исполняется 80~лет со дня рождения академика


И.\,А.~Мизина, директора Института проблем информатики РАН в~период 1989--1999~гг.\ Игорь Александрович ушел из жизни рано, в~возрасте всего 64~лет, но внес огромный вклад как в~развитие теории


отечественных информационно-те\-ле\-ком\-му\-ни\-ка\-ци\-он\-ных сис\-тем


(ИТС), так


и~в~практику их разработки и~внедрения.





\pagebreak





\





\vspace*{-12pt}





\begin{floatingfigure}{68mm}


 \begin{center}


 \mbox{%


 \epsfxsize=63mm


 \epsfbox{miz-1.eps}


 }


 \end{center}


 \vspace*{9pt}


 \end{floatingfigure}








     И.\,А.~Мизин родился 12~апреля 1935~г.\ в~Москве в~семье


военнослужащего. В~1952~г.\ с~отличием окончил среднюю школу и~поступил в~одно из самых престижных военных учебных заведений~---


Во\-ен\-но-воз\-душ\-ную академию име\-ни Н.\,Е.~Жуковского. После


окончания академии в~1959~г.\ по специальности


 <<Эксплуатация


радиотехнических\linebreak средств ВВС>> И.\,А.~Мизин моло\-дым лейтенантом


был направлен для прохождения службы в~одну из ведущих


организаций во\-ен\-но-про\-мыш\-лен\-но\-го комплекса~--- НИИ\linebreak


автоматической аппаратуры, в~ко\-тором


 проработал более 30~лет,


последовательно пройдя ступени карь\-ерного роста от инженера до


замести\-теля директора, руководителя крупного


     на\-уч\-но-про\-из\-вод\-ст\-вен\-но\-го коллектива, главного


конструктора СОД~[1, 2].





     В 1966~г.\ защитил кандидатскую диссертацию на тему


<<Разработка метода обмена информацией по закрытым телеграфным


каналам связи в~командной системе управления специального


назначения>>, а~в~1972~г.~--- докторскую на тему <<Вопросы


исследования и~разработки информационных сетей территориальных


автоматизированных сис\-тем управления>>. В~1975~г.\ И.\,А.~Мизину


присвоено ученое звание профессора по кафедре


<<Автоматизированные сис\-те\-мы управления>>. В~1984~г.\ он был


избран чле\-ном-кор\-рес\-пон\-ден\-том АН СССР по специальности


<<Вы\-чис\-ли\-тель\-ная техника>>.





     В течение последних 10~лет жизни (1989--1999~гг.)\


И.\,А.~Мизин~--- директор ИПИ РАН, генеральный конструктор работ


по модернизации информационных сетей органов государственной


власти, генеральный конструктор АСУ Вооруженных Сил РФ.


В~1997~г.\ был избран академиком РАН по специальности


<<Вы\-чис\-ли\-тель\-ная техника и~элементная база>>.





     Важно, что на всех этапах жизненного пути И.\,А.~Мизин умело


сочетал активную работу ин\-же\-не\-ра-кон\-ст\-рук\-то\-ра систем


и~комплексов передачи данных в~интересах различных заказчиков


(в~том числе в~качестве главного конструктора) с~глубокими


теоретическими исследованиями в~этой предметной области. Он автор


более 190~научных работ в~открытой печати, в~том числе


12~монографий, 12~изобретений, более ста отчетов по НИР, а~также


огромного числа закрытых проектов по созданию систем и~комплексов


различного назначения. Научные работы И.\,А.~Мизина отличались


глубиной и~оригинальностью решений, сочетанием серьезного


математического аппарата с~инженерной интуицией и~практической


направленностью.





     Наиболее известные работы И.\,А.~Мизина:


     \begin{itemize}


\item Передача информации в~сетях с~коммутацией сообщений.~--- М.:


Связь, 1972, 1977 (2-е изд., перераб. и~доп.). Совместно


с~Уринсоном~Л.\,С., Храмешиным~Г.\,К.


\item Сети коммутации пакетов.~--- М.: Радио и~связь, 1986. Совместно


с~Богатыревым~В.\,А. и~Кулешовым~А.\,П.


\item Протоколы информационно-вычислительных сетей: Справочник~/


Под ред.\ И.\,А.~Мизина, А.\,П.~Кулешова.~--- М.: Радио и~связь, 1990.


\item О концепции создания российской общегосударственной и~региональных интегрированных сетей передачи информации~/\!\!/


Электросвязь, 1993. №\,12.


\item Состояние и~перспективы развития информационных


телекоммуникационных технологий для сферы образования и~науки~/\!\!/


Образование и~информатика: Доклады II Международного конгресса


ЮНЕСКО, 1996.


\item Некоторые проблемы создания единой


инфор\-ма\-ци\-он\-но-те\-ле\-ком\-му\-ни\-ка\-ци\-он\-ной системы


общенационального масштаба как основы информатизации сферы


образования в~России~/\!\!/ Системы и~средства информатики. Вып.~8.


\item Информационные технологии в~образовании: от компьютерной


грамотности к информационной культуре общества.~--- М.: Наука, 1996.


C.~114--124. Совместно с~Киселевым~Э.\,В., Соколовым~И.\,А.,


Шоргиным~С.\,Я.


\item Современное состояние проблематики интегрированных


информационно-те\-ле\-ком\-му\-ни\-ка\-ци\-он\-ных систем и~сетей~/\!\!/ Системы и~средства информатики.~--- М.: Наука, 1999. Вып.~9. C.~11--23.


     \end{itemize}





     В настоящей статье научный и~практический вклад академика


И.\,А.~Мизина рассматривается в~рамках трех направлений его


     на\-уч\-но-прак\-ти\-че\-ской деятель\-ности, связанных


с~разработкой и~внедрением СОД с~пакетной


коммутацией в~интересах АСУ Вооруженных Сил, обоснованием


информационных технологий для создания крупных территориальных


систем передачи данных и~созданием информационных сетей


в~интересах органов государственной власти.











\section{И.\,А. Мизин~--- создатель системы обмена данными


в~интересах автоматизированной системы управления стратегического назначения}





     Становление И.\,А.~Мизина и~как инженера-конструктора, и~как


ученого происходило в~стенах НИИ автоматической аппаратуры (НИИ


АА), ведущего в~стране предприятия в~области автоматизированных


систем управления~[3]. В~первые годы И.\,А.~Мизин занимался


исследованиями и~разработкой методов повышения достоверности


передачи информации по каналам связи различной физической


природы, методов помехоустойчивого кодирования, вероятностных


моделей дискретных каналов связи, которые были реализованы в~нескольких устройствах повышения достоверности для


автоматизированных систем с~искажениями. Его целеустремленность и~работоспособность, а~также достигнутые успехи были отмечены


руководством института: в~1963~г.\ он назначается начальником


лаборатории аппаратуры передачи данных.





    Судьбоносным для И.\,А.~Мизина стал 1967~год. Постановлением


ЦК КПСС и~Совета Министров была задана разработка


автоматизированной командной системы боевого управления, а~НИИ


АА определен головной организацией по созданию этой системы. При


формировании коллектива исполнителей И.\,А.~Мизин руководством


института был назначен Главным конструктором СОД, которая должна была


обеспечить беспрецедентно высокие


ве\-ро\-ят\-ност\-но-вре\-мен\-н$\acute{\mbox{ы}}$е


характеристики информационного обмена


данными между объектами, распределенными по территории всей


страны.





    <<\textit{В то уже далекое время, когда в~1967~году в~НИИ


автоматической аппаратуры была поставлена задача создать


телекоммуникационную СОД для


проектируемой тогда большой территориально распределенной АСУ,


проблема ее проектирования представлялась молодому коллективу


разработчиков во главе с~его руководителем И.\,А.~Мизиным не такой


уж сложной. Однако жизнь быстро опровергла это изначальное


пред\-став\-ле\-ние, поскольку предъявленные к~ней


опе\-ративно-технические требования оказались весьма жесткими как в~час\-ти


научно-теоретической проработки системных вопросов, так и~в плане


ее конкретного сис\-те\-мо\-тех\-ни\-че\-ско\-го проектирования и~внедрения.


Практически поиск приемлемых научных и~системотехнических


решений пришлось начинать почти с~нуля}>>,~--- вспоминает


Г.\,К.~Храмешин, один из ближайших соратников Мизина~[4].





    Проведенные предпроектные исследования показали, что


структура создаваемой СОД должна строиться на новых решениях,


отличных от используемых тогда иерархических структур


с~непосредственными связями управляющих звень\-ев АСУ


с~многочисленными объектами ее нижних уровней управления. Такую


структуру практически невозможно было реализовать из-за слишком


больших объемов технических средств пе\-ре\-да\-чи-при\-ема данных


на управляющих объектах обслуживаемой АСУ, которые стали бы


очень громоздкими и~энергоемкими, а~СОД в~целом оказалась бы


системой с~низкими показателями надежности функционирования.





    В связи с~этим И.\,А.~Мизиным совместно с~Л.\,С.~Уринсоном и~Г.\,К.~Храмешиным был предложен принципиально новый в~то время


вариант структурного построения СОД на основе центров коммутации


пакетов (ЦКП) двух модификаций: главных~--- в~районах размещения


объектов высших звеньев АСУ~--- и~территориальных, размещаемых


по всей территории СССР~[5]. Объекты управ\-ле\-ния, являющиеся


абонентами СОД, подключаются дву\-мя--тре\-мя каналами связи


к~двум--трем ближайшим ЦКП. Такое построение позволяло обеспечить


резкое сокращение необходимого количества каналов связи, высокую


надежность связности всех або\-нен\-тов-объек\-тов, но вместе с~тем


приводило к необходимости решения комплекса сложных вопросов


обеспечения процессов их информационного взаимодействия.





    В рамках такой концепции под непосредственным руководством


И.\,А.~Мизина был проработан и~успешно решен целый ряд новых


научно-теоретических и~системотехнических задач, обеспечивших


выполнение множества требований в~интересах АСУ. Основные из


этих задач~[3,~4]:


    \begin{itemize}


\item обоснование структуры циркулирующих в~СОД сообщений и~системы их адресования;


    \item разработка совокупности (стека) алгоритмов, протоколов и~интерфейсов, обеспечивающих передачу через СОД всех необходимых


видов информации;


    \item создание в~составе СОД подсистемы контроля


функционирования и~управ\-ле\-ния ее работой;


    \item создание специализированных аппаратно-программных


комплексов ЦКП и~ап\-па\-рат\-но-программных


комплексов оконечных средств, устанавливаемых на объектах


управления АСУ~--- абонентах СОД;


    \item  разработка комплекса организационно-технических


мероприятий по защите информации, циркулирующей в~СОД.


    \end{itemize}





    Успешное решение этих задач позволило обеспечить высокие


вероятностно-вре\-мен\-н$\acute{\mbox{ы}}$е и~надежностные характеристики доведения


информации, а~также показатели достоверности передачи сообщений


по каналам связи, в~том числе низкого качества.





    Кроме того, коллективом И.\,А.~Мизина был своевременно


выполнен комплекс работ по вводу СОД в~эксплуатацию, который


включал отладку средств и~комплексов на объектах эксплуатации,


проведение различных видов испытаний (автономных,


предварительных, государственных) про\-грам\-мно-ап\-па\-рат\-ных


средств, комплексов и~опытного участка СОД в~целом, а~также


поддержание функционирования системы в~штатном режиме.





    Сравнительно немногочисленный коллектив инженеров, который


во второй половине 1960-х~годов заложил основы новых


системотехнических решений для создания СОД и~ее технических


комплексов, под руководством И.\,А.~Мизина постоянно рос, повышал


свой на\-уч\-но-тех\-ни\-че\-ский потенциал и~к~моменту постановки


системы на опытную эксплуатацию (1980~г.)\ превратился в~мощнейший коллектив высококвалифицированных опытных ученых,


конструкторов, инженеров и~специалистов.





    В процессе создания СОД коллектив, возглавляемый главным


конструктором И.\,А.~Мизиным, работал в~тесном контакте как с~сотрудниками других подразделений института, так и~с учеными и~специалистами институтов и~заводов внешней кооперации. Основные


из них: НПО <<Красная заря>> (г.~Ленинград), НИИССУ (г.~Москва),


НИИ автоматики (г.~Москва), Ереванский НИИ математических


машин, Пензенский НИИЭИ, Тамбовский завод <<Тамбоваппарат>>,


Рязанский завод <<Красное знамя>>, Горьковский НИИРС, КБ завода


<<Сигнал>> (г.~Кишинев), Электромеханический завод


(г.~Свердловск), монтажные предприятия <<Каскад>> и~др.~[3].





    В своей работе И.\,А.~Мизин и~его коллектив разработчиков


постоянно и~конструктивно взаимодействовали с~Генеральным


заказчиком АСУ (маршал А.\,И.~Белов, генералы К.\,Н.~Трофимов,


А.\,П.~Зименков, полковник А.\,П.~Жуковский и~др.)\


и~с~головными военными НИИ, прежде всего с~16~ЦНИИИ связи и~27~ЦНИИ~[6, 7].





     Создание первой в~СССР телекоммуникационной системы с~коммутацией пакетов для командной АСУ и~ее постановка в~режим


круглосуточной эксплуатации стало крупным


     на\-уч\-но-тех\-ни\-че\-ским достижением отечественной


системотехники. И.\,А.~Мизин, как Главный конструктор, в~1981~г.\


стал лауреатом Ленинской премии СССР, а~ряд его сотрудников


удостоены Государственной премии СССР, награждены орденами и~медалями.





     Новые решения, заложенные в~основу построения СОД,


оказались эффективными и~перспективными, обеспечившими


возможность ее эволюционного развития в~базовую


СОД (БСОД) как основы объединенной СОД.





     Именно поэтому БСОД, которая в~1980~г.\ была поставлена на


опытную эксплуатацию, с~1985~г.\ в~режиме боевого дежурства более


четверти века успешно выполняла задачи по информационному обмену


в интересах сотен территориально удаленных объектов АСУ


стратегического назначения~[8].





     Сегодня можно смело утверждать, что многие


системотехнические решения по построению и~функционированию


СОД, разработанные и~успешно внедренные под


руководством И.\,А.~Мизина, предвосхитили ряд идей, которые в~последующем легли в~основу передовых зарубежных разработок в~области телекоммуникационных сетей и~технологий. Эти результаты


остались неизвестными международной научной общественности в~силу закрытого характера работ НИИ~АА.





\section{Информационные технологии для создания крупных


территориальных телекоммуникационных сетей}





     И.\,А.~Мизин был крупным ученым в~области информатики и~внес большой вклад в~развитие


     ин\-фор\-ма\-ци\-он\-но-вы\-чис\-ли\-тель\-ных и~телекоммуникационных систем не только специального, но и~общегосударственного назначения.





     К наиболее значимым достижениям в~этой области следует


отнести развитие теоретических основ


     ин\-фор\-ма\-ци\-он\-но-те\-ле\-ком\-му\-ни\-ка\-ци\-он\-ных


технологий, разработку информационных технологий построения


крупномасштабных ин\-фор\-ма\-ци\-он\-но-вы\-чис\-ли\-тель\-ных и~коммуникационных сетей, а~также создание и~внедрение комплекса


базовых технических и~программных


     средств~\cite{1-miz, 2-miz, 8-miz}.





     Естественно, что методологической базой для этих достижений


стали теоретические результаты и~системотехнические решения по


созданию СОД с~использованием метода пакетной


коммутации, разработанные в~период работы в~НИИ~АА.





     Большое значение для развития отечественной информатики


имели концептуальные работы И.\,А.~Мизина по состоянию и~перспективам развития телeкоммуникационных технологий. В~своем


докладе на эту тему, сделанном на II~Международном конгрессе


ЮНЕСКО <<Образование и~информатика>>  (1996~г.), И.\,А.~Мизин


определил основные этапы развития телекоммуникационных


технологий следующим образом: телеграфные и~телефонные сети


(докомпьютерная эпоха); передача данных между отдельными


абонентами по выделенным и~коммутируемым каналам с~использованием модемов; сети передачи данных с~коммутацией


пакетов; локальные вычислительные сети (ЛВС); цифровые сети


интегрального обслуживания (ISDN~--- Integrated Services for Digital Natwork)~--- узкополосные, а~затем


широкополосные; высокоскоростные распределенные сети~--- Frame


Relay, SMDS (Switched Multimegabit Data Service), ATM (Asynchronous Transfer Mode);


высокоскоростные локальные сети~--- Fast


Ethernet, FDDI (Fiber Distributed Data Interface), FDDI II; информационные супермагистрали.


%


     При этом в~качестве наиболее перспективных технологий


И.\,А.~Мизин определил распределенные сети на оптоволокне,


беспроводные сети мобильных абонентов, Интернет и~электронную


почту. Эти основополагающие выводы остаются актуальными и~сегодня.





     Особое внимание И.\,А.~Мизин уделял развитию информатики и~информационных технологий с~учетом специфических условий России.


Он отмечал, что создание современной телекоммуникационной


инфраструктуры России является сложной масштабной задачей, для


успешного решения которой необходимо комплексное проведение


работ по трем направлениям: реализация крупномасштабных


общегосударственных проектов; развитие и~поддержка региональных


телекоммуникационных проектов; поддержка деятельности


негосударственных организаций в~этой области.





     И.\,А.~Мизиным была поставлена задача описания класса


крупномасштабных телекоммуникационных систем двойного


применения как полносвязных территориально структурированных


мультисетевых систем общенационального масштаба, совмещающих


функции специальных систем и~систем общего


     пользования~\cite{8-miz}.





     Научные работы И.\,А.~Мизина отличались глубиной и~оригинальностью решений, сочетанием серьезного математического


аппарата с~инженерной интуицией, практической направленностью. Он


всегда уделял большое внимание научной обоснованности


предлагаемых решений, в~частности методам математического


моделирования и~оптимизации ИТС.


По инициативе И.\,А.~Мизина в~1992~г.\ в~ИПИ РАН был создан новый отдел~--- <<Отдел


исследований и~разработки принципов построения


     ин\-фор\-ма\-ци\-он\-но-вы\-чис\-ли\-тель\-ных систем>>, одной


из задач которого стало решение научных задач по разработке методов


и~средств моделирования реальных телекоммуникационных систем.


В~частности, проводились фундаментальные теоретические


исследования, разрабатывались математические и~программные


средства и~модели, их экспериментальная апробация, анализ и~оптимизация характеристик ИТС (процессов доставки информации с~разными дисциплинами обслуживания и~входными потоками,


процессов управления и~комплексной обработки информации).





     В этой области под руководством и~при непосредственном


участии И.\,А.~Мизина получен целый ряд результатов~\cite{8-miz}.


Среди наиболее значимых результатов следует отметить комплекс


математических методов, алгоритмов и~программ для расчета


общесетевых характеристик телекоммуникационных сетей большой


размерности, учитывающих реальные технологии, определяемые


действующими международными стандартами и~рекомендациями. Под


руководством И.\,А.~Мизина коллективом ИПИ РАН в~рамках


сотрудничества с~компанией SOTAS (США) были разработаны


концептуальная, математическая и~программная модели


телекоммуникационной сети, использующей популярную в~1990-е~гг.\


технологию передачи Frame Relay.





     Особо необходимо отметить усилия И.\,А.~Мизина по


практической реализации полученных научных результатов. Так,


в~начале 1990-х~гг.\ в~России начались процессы создания


телекоммуникационных сетей, обеспечивающих ограниченный спектр


услуг и~построенных, как правило, на оборудовании зарубежного


производства. И.\,А.~Мизин получил предложение от


Государственного комитета по науке и~технике (ГКНТ) построить сеть


передачи данных в~интересах администрации Псковской области.


К~моменту поступления такого предложения коллективом ИПИ РАН был


обоснован замысел создания отечественного коммутатора,


реализующего протокол~Х.25 и~его модификации. Этот коммутатор и~стал технической основой региональной сети передачи данных в~Псковской области, которая под руководством И.\,А.~Мизина в~период


1994--1995~гг.\ была разработана и~введена в~действие~\cite{10-miz}.





     Рассматривая вклад академика И.\,А.~Мизина в~решение


стратегических проблем информатизации общества в~России, нельзя не


сказать о~разработанной в~ИПИ РАН под его научным руководством


концепции создания единой


ИТС общенационального масштаба (ЕИТС). Основные положения


этой концепции были опубликованы еще в~1996~г., но многие из них


остаются актуальными и~сегодня.





%\vspace*{-16pt}





\section{И.\,А.~Мизин~--- Генеральный конструктор комплекса


работ по~созданию информационных систем органов


государственной власти}





%\vspace*{-6pt}





     На начало 1990-х~гг.\ пришелся переломный момент в~становлении и~использовании информационных технологий в~России в~различных сферах: промышленности, науке, образовании, а~также в~государственном управлении.





     В сентябре 1994~г.\ вышел указ Президента РФ, поставивший


актуальнейшую задачу модернизации информационных и~телекоммуникационных систем органов государственной власти на


базе использования последних достижений в~этой отрасли. Этим же


указом И.\,А.~Мизин был назначен Генеральным конструктором, а~Институт проблем информатики РАН определен головной


организацией.





     Предстояла не только сложная, но и~во многом пионерская


работа~\cite{8-miz}. Необходимо было с~новых позиций изучить


процессы управления, вычленить основные функциональные задачи,


сформулировать количественные и~качественные характеристики,


необходимый набор исходной информации и, наконец, создать


комплексы ап\-па\-рат\-но-про\-грам\-мных средств (в~том числе


сетевого исполнения), позволяющие с~требуемыми характеристиками


автоматизировать эти процессы. В~итоге требовалось создать


компьютерные системы поддержки принятия решений и~автоматизации


процессов государственного управления.





     С одной стороны, требовалось обеспечить преемственность в~использовании уже имеющегося задела, с~другой~--- внедрить самые


современные и~перспективные технологии и~средства для его


реализации. Масштаб и~сложность такой задачи многим


представлялись откровенно пугающими; нередки были голоса,


предрекающие полный крах замысла и~скорое прекращение работы.





     И здесь в~полной мере проявились такие качества личности


И.\,А.~Мизина, как многоплановость, широта и~целеустремленность.


В~нем органично сочетались качества ученого, способного воспринять


и правильно оценить новое, и~инженера, точно знающего, что


необходимо для достижения цели.





     Именно в~процессе этой работы были сформулированы


фундаментальные принципы автоматизации процессов


государственного управления (применительно к России), разработаны


сис\-тем\-но-тех\-ни\-че\-ские решения (вплоть до конкретных


     ап\-па\-рат\-но-про\-грам\-мных решений) построения всех видов


автоматизированных систем, которые долгое время (вплоть до наших


дней) являются категориями первого выбора для многих разработчиков


в нашей стране. Кратко об основных из этих решений~\cite{8-miz}.





     В качестве первого из них следует назвать выбор архитектуры


объектовых комплексов средств автоматизации (КСА), или, в~терминологии, предложенной И.\,А.~Мизиным, комплексов средств


информатизации (КСИ). Что предпочесть~--- централизованную


архитектуру, основанную на одном мощном центральном вычислителе


с подключенными к нему терминалами, или децентрализованную


     кли\-ент-сер\-вер\-ную архитектуру на базе интеллектуальных


автоматизированных рабочих мест (АРМ) и~целевых серверов,


объединенных ЛВС?





     В результате углубленных аналитических исследований и~практических экспериментов, включая анализ опыта крупнейших


зарубежных компаний, в~качестве решения была обоснована


кли\-ент-сер\-вер\-ная архитектура ЛВС. Отказ от многолетней и~широко


распространенной в~Министерстве обороны и~других силовых


ведомствах практики использования терминальных сетей на базе


универсальных ЭВМ (ЕС ЭВМ, IBM) был обусловлен тем, что


произошел коренной перелом в~использовании средств


вычислительной техники и~телекоммуникаций. Существо этого


перелома заключалось, во-пер\-вых, в~том, что подавляющему


большинству пользователей стали доступны мощные персональные


ЭВМ (ПЭВМ), а~во-вто\-рых, в~том, что эти средства перестали


служить только для решения узкоспециализированных задач


высококвалифицированным персоналом. Поэтому при создании


названных выше систем задача состояла в~обеспечении пользователя


(специалиста) всем набором программных средств, позволяющих ему


заниматься своей профессиональной деятельностью. Краткие сроки


создания системы не позволяли начинать все разработки


функционального программного обеспечения <<с нуля>>, а~требовали


применения (адаптации) уже готового программного обеспечения


независимых производителей. Это, в~свою очередь, означало, что перед


создателями системы ставилась не только задача интеграции на одном


АРМ системного программного обеспечения различных


производителей для его совместной работы, но и~задача поддержки


последующего расширения перечня устанавливаемых программных


средств. Имевшиеся на текущий момент тенденции развития


архитектур программных сис\-тем однозначно указывали на


необходимость выбора архитектуры приложения, в~которой


прикладные и~пользовательские сервисы реализованы на клиентской


рабочей станции, а~данные централизованно хранятся и~обрабатываются на сервере.


{\looseness=1





}





     Другим важным вопросом был вопрос выбора технологии


организации ЛВС.





     В отличие от сегодняшней ситуации, когда всецело и~безраздельно


     <<господствует>> протокол Ethernet, в~начале 1990-х~гг.\ на


равных применялись и~Ethernet (тонкий и~толстый), и~Token Ring, и~FDDI, и~ряд других протоколов.





     На основе анализа отечественной и~зарубежной


     на\-уч\-но-тех\-ни\-че\-ской литературы по данной проблематике,


проведения математического моделирования и~стендового


макетирования в~качестве наиболее рационального решения была


выбрана технология Ethernet, в~дальнейшем с~использованием


коммутаторов ЛВС.





     Вместе с~тем стало понятно, что эффективность решения


поставленных задач во многом будет определяться не только


архитектурой и~технологией ЛВС, но и~организацией хранения


информации, методами ее поиска и~извлечения.





     Жесткие ограничения по производительности процессоров


начала 1990-х~гг., ограниченные объемы оперативной памяти при


необходимости реализации сложнейших алгоритмов обработки данных


часто приводили коллективы разработчиков к собственным решениям


по способам структуризации данных в~рамках сетевой модели,


построению уникальных индексных систем, собственных языков


манипулирования данными, их интерпретаторов и~т.\,п. Особенно


остро эти проблемы стояли при обработке неформализованной


текстовой информации на русском языке.





     В этих условиях, несмотря на серьезные возражения, были


обоснованы стратегические решения по унификации подходов к


хранению данных в~рамках реляционной модели и~использованию


языка SQL (Structured Query Language)


как унифицированного средства взаимодействия с~системами управления базами данных (СУБД). В~результате этой


работы при наличии нескольких десятков вариантов в~1995~г.\


в~качестве основных были выбраны три линейки СУБД: СУБД Oracle


(один из лидеров рынка на тот момент) для мощных серверов,


обрабатывающих большие объемы фактографической и~документальной информации с~интенсивным потоком транзакций;


персональная СУБД Microsoft Access как альтернатива


многочисленным линейкам СУБД форматов dBASE и~Paradox для


создания баз данных для небольших групп пользователей; СУБД Microsoft SQL


Server для технологических и~учетных баз данных.





     Впоследствии на основе этих СУБД было разработано множество


приложений, которые, с~небольшими изменениями, работают и~поныне. И~сейчас, спустя почти 20~лет, эти линейки СУБД остаются


доминирующими в~своих сегментах, а~их многочисленные конкуренты


(Sybase SQL, Watcom SQL, Informix, Interbase, Paradox, FoxPro)


практически исчезли с~российского IT-рынка.





     Третьим важнейшим результатом стали решения по контролю


состояния систем, управлению их функционированием и~эксплуатацией. Создаваемые крупномасштабные системы требовали


новых, современных решений в~этой части. И.\,А.~Мизин


фундаментально подошел к решению этих вопросов. Здесь необходимо


отметить, что особое внимание к этим вопросам было обусловлено


опытом Игоря Александровича, приобретенным в~период создания под


его руководством СОД. Разработанная в~ее


рамках система контроля состояния и~управления сыграла огромную


роль в~обеспечении длительного непрерывного и~эффективного


функционирования СОД.





     В 1990-е гг.\ появился целый ряд новых продуктов нового


класса~--- сетевых платформ управления (к~тому времени уже были


специфицированы и~приняты в~качестве международных стандартов


протоколы SNMP~--- Simple Network Management Protocol)).


Наиболее известными были продукты фирмы


Hewlett Packard (HP)~--- линейка продуктов HP Open View~--- и~Computer Associates (CA)~--- программный комплекс UniCenter. ИПИ


РАН были приобретены оба продукта и~проведены всесторонние


аналитические исследования на различных конфигурациях серверного


оборудования и~средствах телекоммуникаций, операционных системах,


СУБД, разнообразном общесистемном и~прикладном программном


обеспечении. Более того, специалисты института прошли


обучение на курсах фирм HP и~СА.





     В результате было принято решение разрабатывать собственную


систему управления с~использованием отдельных модулей из линейки


HP Open View (в~час\-ти управления аппаратными средствами и~контроля трафика) и~модулей из состава продукта UniCenter в~части


управления программным обес\-пе\-че\-нием.


{\looseness=1





}





     Еще одно из важнейших направлений работ связано с~решениями


по обеспечению сопряжения систем, создаваемых на базе современных


технологий, и~систем, уже созданных и~функционирующих на


территории РФ, с~учетом совместного использования их потенциала


(на современном языке~--- сохранения ин\-вес\-тиций).


{\looseness=1





}





     Здесь необходимо отметить два важнейших технических


решения.





     Первое из них~--- сопряжение с~наиболее массовыми и~территориально распространенными сетями того времени:


телеграфными сетями АТ-50 и~Телекс. Такое решение позволяло


повысить вероятность доведения сообщений ограниченного объема


(команды, сигналы, распоряжения и~т.\,п.)\ с~учетом увеличения


множества возможных путей передачи информации за счет


разветвленности телеграфной сети. С~учетом таких потенциальных


возможностей этих сетей в~ИПИ РАН был разработан специальный


комплекс сопряжения с~сетями абонентского телеграфирования АТ-50


и Телекс для обеспечения обмена телеграфными сообщениями между


абонентами объектовых ЛВС (Ethernet) и~абонентами телеграфных


сетей АТ-50 и~Телекс.





     Вторая работа~--- сопряжение с~объектами управления,


подключенными к базовой СОД Вооруженных Сил РФ, которая к


началу 1990-х~гг.\ была основной транспортной сетью для


автоматизированных систем управления Генерального штаба,


различных видов вооруженных сил и~родов войск (АСУ ВС РФ). Сопряжение


обеспечивало возможность информационного обмена с~объектами


военных систем. Для реализации такого сопряжения потребовалась


разработка специального комплекса и~правил по преобразованию


принятых в~базовой СОД форматов кодограмм в~форматы протокола


Ethernet, что позволило обеспечить информационное взаимодействие


абонентов сети Ethernet модернизируемой системы органов


государственной власти и~абонентов КСА автоматизированных систем


военного назначения. И~такой, весьма непростой, комплекс


сопряжения был создан.





     В процессе проектирования было рассмотрено несколько


вариантов реализации этого сопряжения, и~в~качестве основного был


выбран вариант, при котором комплекс сопряжения модернизируемой


системы органов государственной власти подключался по выделенным


каналам связи к центрам коммутации базовой СОД в~виде


виртуального объекта управления, обладающего необходимыми


правами и~характеристиками.





     Естественно, что успешному решению этой задачи


способствовало то, что И.\,А.~Мизин как Главный конструктор СОД


прекрасно владел существом и~особенностями технических решений,


реализованных в~базовой СОД АСУ ВС~РФ.


{\looseness=1





}





     Совершенно новые задачи коллективу ИПИ РАН под


руководством И.\,А.~Мизина пришлось решать в~области


информационной безопасности. Одновременно с~формированием


нормативной базы (совместно с~ФАПСИ, ФСБ и~Министерством


обороны) были апробированы различные подходы и~решения, в~том


числе и~не нашедшие дальнейшего применения. В конечном итоге


вышли на использование решений, которые определяют современное


состояние защищенности информации, содержащей государственную


тайну, в~автоматизированных системах. Разработанные средства


позволяют обеспечить защиту информации при работе пользователей в~доменной структуре, обмене сообщениями электронной почты, доступе


к информации в~базах данных и~при совместной работе над


документами. Средства криптографической защиты реализуют


шифрование и~имитозащиту файлов, функции электронного замка,


шифрование внутрисетевого трафика и~защиту внешних каналов связи.


Впервые в~изделиях были реализованы криптоалгоритмы семейства


<<Метель>>, разработаны специальные считывающие устройства, в~качестве носителей ключевой информации применена Российская


интеллектуальная карта, реализованы различные типы интерфейсов,


включая высокоскоростные оптоволоконные. Разработанные изделия


успешно прошли сертификацию в~ФСБ и~широко применяются в~защищенных сетях органов государственной власти РФ.





     Оперативность и~правильность принятия решений руководством


страны во многом зависит от наглядности и~качества представления


данных, что потребовало широкого внедрения технологий


геоинформационных систем (ГИС), поз\-во\-ля\-ющих эффективно


взаимоувязать картографическую основу и~фактографические базы


данных, а~также быстро отобразить огромный объем накопленных


данных в~пригодном для восприятия виде. Эффективность ГИС в~процессе оперативного принятия решений была настолько очевидной,


что под них создавались специализированные средства отображения


информации~--- <<Видеостены>>, которые позволяли получать


приемлемые характеристики по разрешению, яркости и~размеру


картографического изображения при достаточно скромных


характеристиках единичного технического элемента (проектора,


видеокуба). В~состав интегрированного


     про\-грам\-мно-ап\-па\-рат\-но\-го комплекса, реализующего


функции представления информации для принятия решения, входил


экран коллективного пользования проекционного типа на базе


видеокубов, индивидуальные средства отображения


     и~ап\-па\-рат\-но-про\-грам\-мные средства поддержки.





     Под руководством И.\,А.~Мизина применительно к задачам


государственного управления была показана важность создания


ситуационных центров (СЦ). Ситуационный центр должен представлять собой комплекс,


состоящий из АРМ для руководства и~групп специалистов, а~также из комплексов специального


программного обеспечения, которые обеспечивают оперативный


анализ ситуации, <<проигрывание>> различных сценариев. Главная


задача СЦ~--- поддержание информационного хранилища данных


в~актуальном состоянии на основе сбора информации из различных


источников и~создание условий для интеллектуальной деятельности


специалистов по информационной поддержке принятия


управленческих решений


на основе современных информационных технологий.





     Как видим, И.\,А.~Мизиным были предугаданы и~сформулированы основные принципы и~технические решения для


широчайшего класса автоматизированных систем органов


государственной власти РФ.





\section{Заключение}





     Вся научно-практическая деятельность академика И.\,А.~Мизина


на протяжении более 40~лет была связана с~проблематикой


информационных и~телекоммуникационных сетей. Под его


руководством и~при его непосредственном участии были решены


многие проблемы информатизации процессов государственного


управления, разработаны концептуальные документы по стратегии


взаимосогласованного развития систем и~комплексов связи и~автоматизированного управления войсками, силами и~оружием с~применением новых информационных технологий, выработаны


организационные и~системотехнические принципы координации работ


по автоматизации процессов управления во всех звеньях Вооруженных


Сил РФ, выполнен крупный комплекс работ по созданию базовых


функ\-ци\-о\-наль\-но-ори\-ен\-ти\-ро\-ван\-ных межвидовых


унифицированных ап\-па\-рат\-но-про\-грам\-мных средств


вычислительной техники для перспективных телекоммуникационных


комплексов нового поколения, разработаны алгоритмы автоматической


обработки и~селекции разнородных потоков текстовой информации и~моделирования критических ситуаций, разнесенных во времени и~пространстве, созданы математические модели сложных адаптивных


телекоммуникационных сетей.


{\looseness=1





}





     И.\,А.~Мизин постоянно вел большую


     на\-уч\-но-ор\-га\-ни\-за\-ци\-он\-ную работу. Он был


председателем секции <<Теория передачи и~обработки информации>>


Научного совета по комплексной проблеме <<Кибернетика>> АН


СССР, членом комиссии по вычислительным центрам коллективного


пользования и~сетям ЭВМ при Президиуме АН СССР, членом Совета


руководителей Академсети, членом советской рабочей группы


международной организации по стандартизации, председателем


Научного совета РАН по информационно-вычислительным сетям,


членом Президиума на\-уч\-но-тех\-ни\-че\-ско\-го совета по программе


<<Информатизация России>>, членом Координационного совета по


информатизации при Администрации Президента РФ, членом научных


советов при Совете Безопасности РФ, Совете обороны РФ,


Минэкономики РФ, членом Президиума НТС при Правительстве


г.~Москвы.





     С 1996~г.\ и~до последних дней И.\,А.~Мизин вел активную


работу как член бюро Отделения информатики, вычислительной


техники и~автоматизации РАН.





     И.\,А.~Мизин много внимания уделял подготовке научных и~инженерных кад\-ров. Он был заведующим базовыми кафедрами в~Московском государственном университете радиотехники,


электроники и~автоматики (МИРЭА) и~в~Московском техническом


университете связи и~информатики (МТУСИ), председателем двух


специализированных советов по защите докторских диссертаций.


Подготовил значительное число кандидатов и~докторов технических


наук.





     И.\,А.~Мизин обладал удивительной способностью создавать в~науке и~промышленности успешно работающие большие коллективы


специалистов, объединять вокруг себя талантливую молодежь.





     После прихода в~ИПИ РАН И.\,А.~Мизин установил и~поддерживал активные контакты с~зарубежными коллегами. Он


получил международную известность и~снискал глубокое уважение


многих ученых, профессоров и~деловых партнеров в~США, Германии,


Италии, странах Юго-Вос\-точ\-ной Азии.





     Научная, научно-организационная и~педагогическая деятельность


И.\,А.~Мизина была по достоинству оценена нашей Родиной. Он был


лауреатом Ленинской и~Государственной премий, награжден орденом


Трудового Красного Знамени, орденом <<За заслуги перед


Отечеством>> 4-й степени, многими медалями.





     Игорь Александрович Мизин умер в~расцвете творческих сил


8~сентября 1999~г.\ после тяжелой болезни. Он похоронен в~Москве на


Востряковском кладбище.





     Коллеги и~ученики И.\,А.~Мизина продолжают начатые им


работы, реализуя и~развивая его научные идеи в~различных


организациях Российской академии наук, промышленности


и~в~технических университетах.





{\small\frenchspacing


{%\baselineskip=10.8pt


\addcontentsline{toc}{section}{Литература}


\begin{thebibliography}{99}





\bibitem{1-miz}


\Au{Захаров В.\,Н., Соколов И.\,А.} Академик Игорь Александрович


Мизин~/\!\!/ История науки и~техники, 2008. №\,7. С.~53--56.


\bibitem{2-miz}


\Au{Зацаринный А.\,А., Захаров В.\,Н., Синицин~И.\,Н.} ИПИ РАН и~создание ин\-фор\-ма\-ци\-он\-но-те\-ле\-ком\-му\-ни\-ка\-ци\-он\-ных


сетей в~интересах органов государственной власти Российской


Федерации (к~25-ле\-тию ИПИ РАН)~/\!\!/ Ведомственные


корпоративные сети, 2008. №\,2(47). С.~118--128.


\bibitem{3-miz}


Автоматизация управления. Наш путь: К~50-ле\-тию НИИ~автоматической аппаратуры


им.~академика В.\,С.~Семенихина.~--- М.: НИИ~АА, 2006. 210~с.


\bibitem{4-miz}


\Au{Храмешин Г.\,К.} Мизин~--- главный конструктор первой в~СССР


телекоммуникационной системы обмена данными на принципах


коммутации пакетов. Основные научные и~системотехнические


проблемы ее создания~/\!\!/ И.\,А.~Мизин~--- ученый, конструктор,


человек.~--- М.: ИПИ РАН, 2010. С.~129--132.


\bibitem{5-miz}


\Au{Щукин Л.\,Б.} И.\,А.~Мизин: системное мышление, новаторство,


целеустремленность~/\!\!/ И.\,А.~Мизин~--- ученый, конструктор,


человек.~--- М.: ИПИ РАН, 2010. С.~136--141.





\bibitem{7-miz} %6


\Au{Зацаринный А.\,А.} О роли и~вкладе института в~создание и~развитие системы обмена данными АСУ ВС~/\!\!/ Институт военной


связи. История и~современность. 1923--2008~гг.~--- Мытищи, 2008.


С.~230--237.





\bibitem{6-miz} %7


\Au{Зацаринный А.\,А.} Академик И.\,А.~Мизин: военная наука и~практика~/\!\!/ И.\,А.~Мизин~--- ученый, конструктор, человек.~--- М.:


ИПИ РАН, 2010. С.~96--128.








\bibitem{9-miz} %8


\Au{Зацаринный А.\,А.} Аппаратура систем и~комплексов обмена


данными~/\!\!/ История отечественных средств связи.~--- М.:


Столичная энциклопедия, 2013. С.~501--506.


\bibitem{8-miz} %9


\Au{Соколов И.\,А.} Вклад академика И.\,А.~Мизина в~развитие


отечественных информационных технологий и~систем~/\!\!/


И.\,А.~Мизин~--- ученый, конструктор, человек.~--- М.: ИПИ РАН,


2010. С.~85--95.





\bibitem{10-miz}


\Au{Гайкович Ю.\,В.} Региональные телекоммуникации от Мизина~/\!\!/


И.\,А.~Мизин~--- ученый, конструктор, человек.~--- М.: ИПИ РАН,


2010. С.~193--196.





\end{thebibliography}


} }








\vspace*{-3pt}





\hfill{\small\textit{Поступила в~редакцию 30.01.15}}








\newpage





%\vspace*{9pt}





%\hrule





%\vspace*{2pt}





%\hrule





\vspace*{-24pt}








\def\tit{ABOUT ACADEMICIAN I.\,A.~MIZIN'S CONTRIBUTION TO~THEORY


AND~PRACTICE OF~DOMESTIC INFORMATION-TELECOMMUNICATION SYSTEMS CREATION:\\


TO~THE~80th ANNIVERSARY}





\def\titkol{About academician I.\,A.~Mizin's contribution to theory and


practice of domestic ITS} %information-telecommunication systems creation: To the 80th anniversary}








\def\aut{I.\,A.~Sokolov, A.\,A.~Zatsarinny, and~V.\,N.~Zakharov}


\def\autkol{I.\,A.~Sokolov, A.\,A.~Zatsarinny, and~V.\,N.~Zakharov}








\titel{\tit}{\aut}{\autkol}{\titkol}





\vspace*{-6pt}





\noindent


Federal Research Center ``Computer Science and Control,'' Russian


Academy of Sciences, 44-2 Vavilov Str., Moscow 119333, Russian


Federation











\def\leftfootline{\small{\textbf{\thepage}


\hfill Sistemy i Sredstva Informatiki~---


Systems and Means of Informatics 2015 vol~25 no\,1}


}%


 \def\rightfootline{\small{Sistemy i Sredstva Informatiki~---


 Systems and Means of Informatics   2015   vol~25   no\,1


\hfill \textbf{\thepage}}}








\Abste{The article is devoted to the 80th anniversary of academician I.\,A.~Mizin,


head of


IPI RAS during 1989--1999, outstanding scientist, designer and engineer. A~brief


biography concerning his scientific work is presented. The scientific and practical


contribution of I.\,A.~Mizin to the theory and its applications in creation of domestic


information-telecommunication systems (ITS)


is considered in three directions of his work. The


first direction is development and implementation of the data communication system for


the purposes of Armed Force ACS (automated control system),


which was the first domestic network with packet


commutation. The second direction is justification of information technologies for


creation of the huge territorial system of data communication in


the regions of Russia. The


third direction is creation and development of information networks for the purposes of


government with requirements of information protection.}





\KWE{chief designer; academician; information telecommunication networks;


information technologies; data communication system; packet commutation; methods of


data communication; information protection; link channels; trial-run; government}





\DOI{10.14357\!/\!08696527150114}





\vspace*{-12pt}





%\Ack


%\noindent








\renewcommand{\bibname}{\protect\rmfamily References}


%\renewcommand{\bibname}{\large\protect\rm References}





{\small\frenchspacing


{%\baselineskip=10.8pt


\addcontentsline{toc}{section}{References}


\begin{thebibliography}{99}





\bibitem{1-miz-1}


\Aue{Zakharov, V.\,N., and I.\,A.~Sokolov}. 2008. Akademik Igor'


Aleksandrovich Mizin [Academician Igor Alexandrovich Mizin].


\textit{Istoriya Nauki i Tekhniki} [Science and Technics History] 7:53--56.


\bibitem{2-miz-1}


\Aue{Zatsarinny, A.\,A., V.\,N.~Zakharov, and I.\,N.~Sinitsyn}. 2008. IPI


RAN i~sozdanie informatsionno-telekommunikatsionnykh setey


v~interesakh organov gosudarstvennoy vlasti Rossiyskoy Federatsii


(k~25-le\-tiyu IPI RAN) [IPI RAS and creation of information


telecommunication networks for Russian Government purposes (to the 25th IPI


RAS anniversary)]. \textit{Vedomstvennye Korporativnye Seti}


[Departmental Corporate Networks] 2(47):118--128.


\bibitem{3-miz-1}


2006. \textit{Avtomatizatsiya upravleniya. Nash put'. K~50-le\-tiyu NII avtomaticheskoy


apparatury


im.~akademika V.\,S.~Semenikhina} [Control automation. Our path. To the 50th


anniversary of academician V.\,S.~Semenikhin


Research Institute of Automated Apparature].


Moscow: NII~AA. 210~p.


\bibitem{4-miz-1}


\Aue{Khrameshin, G.\,K.} 2010. Mizin~--- glavnyy konstruktor pervoy


v~SSSR telekommunikatsionnoy sistemy obmena dannymi na printsipakh


kommutatsii paketov. Osnovnye nauchnye i sistemotekhnicheskie problemy


ee sozdaniya [Mizin as a~chief designer of the first in USSR


telecommunication system of data exchange made by principles of packet


commutation. Main science and systemotechnic problems of its creation].


\textit{I.\,A.~Mizin~--- uchenyy, konstruktor, chelovek} [Mizin as a~scientist, a~designer, a~man]. Moscow: IPI RAS. 129--132.


\bibitem{5-miz-1}


\Aue{Shchukin, L.\,B.} 2010. I.\,A.~Mizin: Sistemnoe myshlenie,


novatorstvo, tseleustremlennost' [I.A. Mizin: Systemic thinking, innovation,


sense of purpose].  \textit{I.\,A.~Mizin~--- uchenyy, konstruktor, chelovek}


[Mizin as a~scientist, a~designer, a~man]. Moscow: IPI RAS. 136--141.








\bibitem{7-miz-1} %6


\Aue{Zatsarinny, A.\,A.} 2008. O~roli i~vklade instituta v sozdanie


i~razvitie sistemy obmena dannymi ASU VS [About institute role and


contribution to creation and development of data exchange system of Armed


Forces ACS].


\textit{Nauchno-istoricheskiy sbornik <<Institut voennoy svyazi. Istoriya


i~sovremennost'. 1923--2008~gg.} [Science-historical digest ``Institute of


military communication. History and modernity'']. Mytishchi. 230--237.





\bibitem{6-miz-1} %7


\Aue{Zatsarinny, A.\,A.} 2010. Akademik I.\,A.~Mizin: Voennaya nauka


i~praktika [An academician I.\,A.~Mizin: Military science and practice].


\textit{I.\,A.~Mizin~--- uchenyy, konstruktor, chelovek} [Mizin as a~scientist, a~designer, a~man]. Moscow: IPI RAS. 96--128.





\bibitem{9-miz-1} %8


\Aue{Zatsarinny, A.\,A.} 2013. Apparatura sistem i~kompleksov obmena


dannymi [Apparature of data exchange systems and complexes].


\textit{Istoriya otechestvennykh sredstv svyazi}


[The history of domestic communication means]. Moscow:


Stolichnaya entsiclopediya. 501--506.


\bibitem{8-miz-1} %9


\Aue{Sokolov, I.\,A.} 2010. Vklad akademika I.\,A.~Mizina v~razvitie


otechestvennykh informa\-tsi\-on\-nykh tekhnologiy i~sistem


[Academician I.\,A.~Mizin's contribution to development of domestic


information technologies and systems]. \textit{I.\,A.~Mizin~--- uchenyy,


konstruktor, chelovek} [Mizin as a~scientist, a~designer, a~man]. Moscow:


IPI RAS. 85--95.





\bibitem{10-miz-1}


\Aue{Gaykovich, Yu.\,V.} 2010. Regional`nye telekommunikatsii ot Mizina


[Regional telecommunications by Mizin]. \textit{I.\,A.~Mizin~--- uchenyy,


konstruktor, chelovek} [Mizin as a~scientist, a~designer, a~man]. Moscow:


IPI RAS. 193--196.


\end{thebibliography}


} }








\vspace*{-6pt}





\hfill{\small\textit{Received January 30, 2015}}





\vspace*{-12pt}





\Contr





\noindent


\textbf{Sokolov Igor A.} (b.\ 1954)~--- Academician of the Russian Academy of Sciences, Doctor


of Science in technology, Director, Federal Research Center ``Computer Science


and


Control,'' Russian Academy of Sciences, 44-2 Vavilov Str., Moscow


119333, Russian Federation; isokolov@ipiran.ru





\vspace*{3pt}





\noindent


\textbf{Zatsarinny Alexander A.} (b.\ 1951)~--- Doctor of Science in technology, professor,


Deputy Director, Federal Research Center ``Computer Science and Control,''


Russian Academy of Sciences, 44-2 Vavilov Str., Moscow 119333, Russian


Federation; azatsarinny@ipiran.ru








\vspace*{3pt}





\noindent


\textbf{Zakharov Victor N.} (b.\ 1948)~--- Doctor of Science in technology, associate professor;


Scientific Secretary, Federal Research Center ``Computer Science and Control,''


Russian Academy of Sciences, 44-2 Vavilov Str., Moscow 119333, Russian


Federation; vzakharov@ipiran.ru





 \label{end\stat}








\renewcommand{\bibname}{\protect\rm Литература}











